\section{Consequences and further questions}\label{sec:consequences}

The counting argument also applies when the algorithm starts or finishes
with a partial order. The two settings have different information bounds:
initial relations reduce what remains to be learned, whereas a target
partial order specifies what the algorithm must establish. We state the
consequences here and give their proofs in Appendix~\ref{app:consequences}.

\subsection{Producing a prescribed partial order}

Let $R$ be a partial order on $n$ slots. In \emph{partial-order production},
the input consists of $n$ keys with no known relations. The algorithm must
assign one key to each slot so that the key in slot $u$ is smaller than
the key in slot $v$ whenever $u<_R v$. Sorting is the case where $R$ is
a chain; a slot required to precede every other slot asks for the minimum.

A fixed assignment is valid for exactly $e(R)$ total orders of the keys.
The information-theoretic lower bound is therefore
\[
I(R)=\log(n!)-\log e(R).
\]
Cardinal et al.~\cite{cardinal-production} give a deterministic
polynomial-time algorithm using $I(R)+o(I(R))+O(n)$ comparisons.

\begin{corollary}[Partial-order production]\label{cor:production}
There is a nonnegative sequence $\varepsilon_n\to0$, independent of the
target, such that the adversary of Theorem~\ref{thm:main} forces every
correct comparison algorithm producing an $n$-element partial order $R$
to make at least
\[
(1-\varepsilon_n)I(R)
\]
comparisons. The adversary is initialized using only $n$ and works for
every target $R$.
\end{corollary}

Its leading term matches Cardinal et al.'s upper bound when
$I(R)=\omega(n)$. The information bound can be smaller than the task's
comparison complexity. For example, producing a minimum has
$I(R)=\log n$ and requires $n-1$ comparisons.

\subsection{Sorting from supplied partial information}

Now let $P$ describe relations already known between the input keys,
supplied as a DAG. The sorter must identify one of its $e(P)$ linear
extensions, so the information bound is $\log e(P)$.

\begin{corollary}[Supplied partial information]\label{cor:initial-poset}
There is a deterministic adversary that can be initialized from any
$n$-element partial order $P$ in polynomial time, gives legal answers in
polynomial time per comparison, and forces every correct sorter to make
at least
\[
(1-o(1))\log e(P)-O(n\log\log n)
\]
additional comparisons. Both error terms are uniform over $P$.
\end{corollary}

The construction uses Cardinal et al.'s polynomial-time procedure~\cite{cardinal-production} to
extend $P$ to a \emph{weak order}: a sequence of antichain layers, with
every key in an earlier layer preceding every key in a later one.
The additional relations reduce $\log e(P)$ by at most
$O(n\log\log n)$.
Run an independent copy of our adversary within each layer and answer
between layers from their fixed order. All copies use the same parameters,
chosen from the total size $n$, so the guarantee also covers small layers.

The comparison bound approaches $\log e(P)$ when
$\log e(P)=\omega(n\log\log n)$. Kahn and Kim's adversary already forces
$\Omega(\log e(P))$ for every $P$. The improvement here is the
coefficient approaching one in this regime.

\subsection{Preserving small amounts of remaining uncertainty}

Can a deterministic polynomial-time adversary force
$(1-o(1))\log e(P)$ further comparisons for every supplied partial order
$P$? The additive loss in Corollary~\ref{cor:initial-poset} can overwhelm
the information that remains. Two disjoint chains of length $n/2$ have
$\log e(P)=\Theta(n)$; a chain on $n-1$ items and one unrelated item have
$e(P)=n$.

Reducing the additive loss to $O(n)$ would already extend the useful
range. Cardinal et al.~\cite[Section~4]{cardinal-production} exhibit
partial orders $P$ for which every weak-order extension $W$ satisfies
$\log e(W)\le\log e(P)-\Omega(n\log\log n)$. A sharper construction must
therefore preserve more of the possible interleavings between layers.

Initial partial information also arises when sorting points whose
locations are known only to lie in specified intervals
\cite{imprecise-points,interval-orders}. Extending our adversary to those
models would require answers compatible with realizable coordinates and
with the permitted geometric operations. Another direction is the
amortized minimum-maintenance problem linked to sorting adversaries by
Kaplan, Zamir, and Zwick~\cite[Section~7]{kzz-min}: handling insertions and
deletions would require controlling how they change the represented orders.
